\documentclass[11pt, a4paper]{amsart}
\usepackage[a4paper,top=2cm,bottom=2cm]{geometry}
\usepackage[backend=biber,autocite=plain,backref=true]{biblatex}
\usepackage[color=yellow,obeyDraft]{todonotes} % white, lime, pink
\usepackage[dvipsnames]{xcolor}
\usepackage[utf8]{inputenc}
\usepackage{amsmath, amssymb, amsfonts}
\usepackage{booktabs} % For enhanced table rules
\usepackage{caption} % un-numbered figures
\usepackage{enumitem}
\usepackage{longtable}
\usepackage{marginnote}
\usepackage{nth}
\usepackage{showframe}
\usepackage{url}

\renewcommand{\thesection}{\arabic{section}}
\renewcommand{\ShowFrameColor}{\color{white}}
\newcommand{\norm}[1]{\left\lVert#1\right\rVert}
\newcommand{\Lap}{\nabla^2}
\newcommand{\Grad}{\nabla}
\newcommand{\del}{\partial}
\DeclareMathOperator{\tr}{tr}
\DeclareMathOperator{\trace}{tr}

%\raggedbottom
\addbibresource{rayleigh.bib}

\title{Primer for Theory of Sound 1}
\author{David checking Gemini}
\date{}

\begin{document}

This document contains miscellaneous text that I find helpful.  It begins with
my notes on Rayleigh's 1881 article on plates and shells.  Backround
information about the historical context follows.  The last section contains my
notes on plates as described in Sound.

\section{Collected Works page 551},
%WHEN \cite{rayleigh1902scientific, rayleigh1881}

With bibtex related,
does this work\cite{Rayleigh1881,ScientificPapersVol1,RayleighSoundVol1Ed2}
%or this\cite{Rayleigh1881,RayleighSoundVol1Ed2}
%or is it\cite{ScientificPapersVol1,RayleighSoundVol1Ed2}

When a thin sheet of matter is subjected to stress, the force which it opposes
to extension is great in comparison with that which it apposes to bending.
Under ordinary circumstances, the deformation takes place approximately as if
the sheet were inextensible as a whole, a condition which in a remarkable
degree facilitates calculation, though (it need scarcely be said even bending
implies an extension of all but the central layers.
The inextensibility postulated refers properly to the central layer, which is
spoken of as a surface.
We will commence with the case of a surface naturally spherical, and
investigate what kind of deformation it admits of, under the condition that no
line traced upon the surface is altered in length. The radius of the sphere
being \(a\), let the point whose natural coordinates are \(a\), \(\theta\), \(\phi\)
be displaced to the position
\[ (a + \delta r, \theta + \delta\theta, \phi + \delta \phi),\]
 where \((\delta r,\delta \theta,\delta \phi)\) is to be treated as small.
\[
M(a,\theta) =
\begin{bmatrix}
a^2 & 0
0  & a^2 \sin^2(\theta)
\end{bmatrix}
\]
Since the element of arc \(ds\) is of the same length after displacement as
\[
\norm{ds}^2 = (d\theta, d \phi) M [d\theta; d \phi ]
\]
I am just making up some notation here;  why not \(x\) ?
\[
x^T = (d\theta + d \delta \theta, \phi + d \delta \phi ),
\]
Inextensibility means
\[
\norm{ds}^2 = x^T  M(a +  \delta r,\theta + \delta\theta) x
\]
Rayleigh includes the addition term \( (d \delta r)^2 \); I do not see
where this comes from; I will just assume that it is \(0\), and see
how far I get.


What is \(x\)?
\[
\Grad = (\frac{\del}{\del \theta}, \frac{\del}{\del \phi}),
\]
\[
df = (\Grad f)  [d \theta; d \phi],
\]
\[
x^T = (d\theta + d \delta \theta, d \phi + d \delta \phi ),
\]
\begin{align}
d\theta + d \delta \theta = ((1,0) +\Grad (\delta \theta)) [d \theta; d \phi] \\
d \phi + d \delta \phi  = ((0,1) + \Grad (\delta \phi)) [d \theta; d \phi]
\end{align}
\[
x = \left( I +
D \begin{matrix}
d\theta &
d \phi
\end{matrix} \right)
\begin{matrix}
d\theta &
d \phi
\end{matrix}
\]
\[
= (\frac{\del f}{\del \theta}, \frac{\del f}{\del \phi}
(\frac{\del f}{\del \theta}, \frac{\del f}{\del \phi}
\]
\[ d f = \frac{\del f}{\del \theta} d \theta  + d f = \frac{\del f}{\del \theta} d \theta \]
\[ d \delta \theta =(\Grad \delta \theta )[d \theta;d \phi]
\frac{ \del \delta \theta }{ \del \theta } d \theta
\]

\begin{verbatim}
= (d\theta, d \phi) M(a +  \delta r,\theta + \delta\theta) [d\theta; d \phi ]
, \phi + d \delta \phi ).
 + d \delta ):,
 + d \delta \theta) + (d \delta r) + (a + \delta r) sin' (\theta + \delta\theta)
(d \phi + d \delta ):,

= (a + \delta r)
(d\theta + d \delta \theta) + (d \delta r) +
(a + \delta r) sin' (\theta + \delta\theta)
(\phi + d \delta \phi ):,
\end{verbatim}
before, we have


This is raw text from an old ocr scanner.
\begin{verbatim}
= (a + \delta r) (d\theta + d \delta \theta) +
(d \delta r) +
(a + \delta r) sin' (\theta + \delta\theta) (2\phi + d \delta ):,

or, retaining only the first power of \delta r, \delta\theta, \delta\phi,


 + ir (de) + ir sin e (d theta) +
a sin \theta cos \theta \delta\theta (d)* +
a sin* \theta dd = \theta.
\end{verbatim}
Now
\[
= (a + \delta r) (d\theta +
d \delta \theta) + (d \delta r) +
(a + \delta r) sin' (\theta + \delta\theta) (2\phi + d \delta ):,
\]
just write e for delta ?

\maketitle

\section{Prologue}\label{prologue}

I found a paper \cite{Woodhouse2024}; frf, matrices, \(\partial\).
Prof. Woodhouse' publications might be more relevant that Rayleigh's book.

\url{https://scholar.google.com/citations?user=2jBVQo4AAAAJ&hl=en}

Also see \url{https://euphonics.org/}



This document collects Gemini responses to prompts discussing prerequisities.

My primary motivation for reading Theory of Sound is to learn the
approximation for the first fundamental mode of structures \cite{Hodges1985}.
The static deflection curve is the deformation of the structure under a uniform
load, e.g., the force of gravity. Rayleigh's Method for estimating the
fundamental frequency of a structure is to use the Rayleigh quotient
corresponding to the static deflection curve.

Another reason to reading Theory of Sound is that it defined the
\emph{lingua franca} for a large proposition of the work related to
acoustics and vibrations done in the early 20th century. Readers must
know Rayleigh's notation `by heart'. Moreover Rayleigh assumed that his
readers were familiar with the famous 17th and 19th century inventors
and their intriguing inventions.

Humans hear sounds in the range 20~Hz to 20,000~Hz. Decibels (dB) are a
measure of wave amplitude, the displacement of particles from their
resting position. Meters \(m\), seconds \(s\), A medium's density and
elasticity determine the speed of sound.

\begin{itemize}
\item
  air (at 20\(^\circ \)C): \(\sim \)343 \(m/s\)
\item
  water: \(\sim \)1,480 \(m/s\)
\item
  steel: \(\sim \)5,960 \(m/s\)
\end{itemize}

As far as I know, heat is the primary cause of variation in sound speed;
temperature is inversely proportional to density.

At a hard or soft boundary, or at the interface between a different
medium, sound waves change at a

\begin{itemize}
\item
  Reflection: Bouncing off a surface, creating echoes or reverberation.
\item
  Refraction: Bending as it passes between zones of different
  temperatures or materials.
\item
  Diffraction: Bending around obstacles or spreading through openings.
\item
  Absorption: Losing energy as heat when hitting soft, porous materials.
\end{itemize}

Diffraction, e.g., Borda's mouthpiece, creates interesting
singularities. The list of intriguing physical phenomena includes
impedance, dispersion, sub-harmonics, Helmholtz's resonator and
thermoacoustics.

\section{Nomenclature}\label{nomenclature}
One motivation for defining Rayleigh's notation is to understand the work of
the other Rayleigh.
\[(D_x u ) + {({D_x u})}^T =
\begin{bmatrix}
2e & c & b\\
c & 2f & a\\
b & a & 2g\\
\end{bmatrix}.\]
\[e=\frac{\partial u}{\partial x},\quad f=\frac{\partial v}{\partial y},\quad
g=\frac{\partial w}{\partial z}\]
\[a=\frac{\partial w}{\partial y}+\frac{\partial v}{\partial z},\quad
b=\frac{\partial u}{\partial z}+\frac{\partial w}{\partial x},\quad
c=\frac{\partial v}{\partial x}+\frac{\partial u}{\partial y}\]
Dilatation (Volumetric Strain): \(\theta =e+f+g\)

Stress (Internal Forces per Unit Area)

To represent the internal forces acting across planes within a vibrating
solid Rayleigh tracked the six components of the stress system using
capital letters:

Normal Stresses: Represented by \(P, Q, R\) (pulling or pushing
perpendicular to the \(x, y, z\) planes).

Tangential (Shear) Stresses: Represented by \(S, T, U\) (acting parallel
to the faces of an infinitesimal cube).

When writing the potential energy of a deformed solid (the Strain Energy
Function, \(W\)), he linked these stresses directly to his strain
components via Hooke's Law using the notation:

\(P=\frac{\partial W}{\partial e},\quad Q=\frac{\partial W}{\partial f},\quad R=\frac{\partial W}{\partial g}\)

\(S=\frac{\partial W}{\partial a},\quad T=\frac{\partial W}{\partial b},\quad U=\frac{\partial W}{\partial c}\)

Body/External Forces: When expressing the equations of motion for an
elastic solid body, Rayleigh denoted the external components of force
acting per unit volume as \(X, Y, Z\).

The Equations of Motion: \(\rho \) is the mass density of the material.
Using his combined notation for stress and force components, Rayleigh
wrote the foundational 3D elasto-dynamic equations as:
\[\rho \ddot{u}=\rho X+
\frac{\partial P}{\partial x}+
\frac{\partial U}{\partial y}+
\frac{\partial T}{\partial z}\]

\begin{longtable}[]{@{}ccc@{}}
\caption{Rayleigh elasticity notation and contemporary equivalents.}\label{tbl:nomenclature1}
\tabularnewline%
quantity & symbol & modern \\
\midrule\noalign{}
\endfirsthead%
\toprule\noalign{}
quantity & symbol & modern \\
\midrule\noalign{}
\endhead%
\endlastfoot{}% chktxt 1
displacement & \(u, v, w\) & \(u_1, u_2, u_3\) \\
Normal Strain & \(e, f, g\) &
\(\epsilon_{xx}, \epsilon_{yy}, \epsilon_{zz}\) \\
Shear Strain & \(a, b, c\) &
\(\gamma_{yz}, \gamma_{xz}, \gamma_{xy}\) \\
Volumetric Strain & \(\theta = \Delta \) &
\(\epsilon_{v}= \text{Tr}(\boldsymbol{\epsilon})\) \\
Normal Stress & \(P, Q, R\) &
\(\sigma_{xx}, \sigma_{yy}, \sigma_{zz}\) \\
Shear Stress & \(S, T, U\) & \(\tau_{yz}, \tau_{xz}, \tau_{xy}\) \\
External Body Force & \(X, Y, Z\) & \(\mathbf{f} = b_{i}\) \\
Acoustic Pressure & \(p\) & \(p\), \(\Delta p\) \\
\end{longtable}


\begin{longtable}[]{@{}cc@{}}
\caption{Notation}\label{tbl:notation}\tabularnewline%
symbol & quantity modern \\
\midrule\noalign{}
\endfirsthead%
\toprule\noalign{}
symbol & quantity modern \\
\midrule\noalign{}
\endhead%
\endlastfoot%
\(T\) & kinetic energy \\
\(V\) & Strain energy \\
\(F\) & dissipation function \\
\(\psi_k\) & generalized coordinate \(k\) \\
\(y\), \(w\) & lateral displacement \\
A & amplitude \(A \cos( p t - \epsilon)\) \\
p & circular frequency \\
\(\epsilon \) & phase angle \\
\(\rho \) & mass density \\
\(\omega \) & cross sectional area \\
E & Young's \\
I & inertia \\
\(\kappa \) & radius of gyration \\
\end{longtable}


When writing the Strain Energy Function (W)---which measures the
potential energy stored in an elastic body when it is deformed.

Rayleigh formulated it using his strain components

(e, f, g for normal strains and a, b, c for shear strains)

alongside m and n:

\(2W=(m+n)(e^{2}+f^{2}+g^{2})+2(m-n)(f g+g e+ef)+n(a^{2}+b^{2}+c^{2})\)

Or in terms of the volumetric dilatation \(\theta = e + f + g\),

\(2W=(m-n)\theta^{2}+2n(e^{2}+f^{2}+g^{2})+n(a^{2}+b^{2}+c^{2})\)

\begin{longtable}[]{@{}ccc@{}}
\caption{Rayleigh's notation (blanks are intentional).}\label{tbl:archaic}
\tabularnewline%
quantity & symbol & modern \\
\midrule\noalign{}
\endfirsthead%
\toprule\noalign{}
quantity & symbol & modern \\
\midrule\noalign{}
\endhead%
\endlastfoot{}%
                 &\(m\)& \(\lambda + \mu \)  \\
shear modulus    &\(n\)& \(\mu = G\) \\
bulk modulus     &     & \(K = \lambda + \frac{2}{3} \mu \) \\
Poisson's ratio  &     & \(\nu = \frac{\lambda }{2 (\lambda + \mu )} \) \\
\end{longtable}

\section{Vocabulary Advance
Organizer}\label{vocabulary-advance-organizer}

\begin{longtable}[]{@{}ccc@{}}
\caption{Inventors and their inventions by chapter.}\label{tbl:nomenclature3}\tabularnewline%
Chapter & & dates, invention \\
\midrule\noalign{}
\endfirsthead%
\toprule\noalign{}
Chapter & & dates, invention \\
\midrule\noalign{}
\endhead%
\endlastfoot%
1 & Regnault & 1810--1878 \\
1 & Cagniard de la Tour & 1819, siren \\
2 & Wheatstone & 1827, kaleidophone \\
3 & Scheibler & 1834, tonometer \\
6 & Mersenne & 1588--1648, sonometer \\
7 & Terquem & 1831--1887 Savart's son rauque \\
9 & Bourget & 1822--1887 \\
10 & Chladni, Mathieu & \\
10 & Rudolph Koenig & 1832--1901 \\
\end{longtable}

Henri Victor Regnault (1810--1878) was a highly precise French chemist
and physicist, famous for his meticulous empirical measurements of the
physical properties of gases.

\textbf{His Contribution}: Regnault conducted massive, definitive
experiments on the velocity of sound in the atmosphere and inside long
pipes (famously utilizing the newly constructed Paris water mains). He
discovered that the speed of sound is not perfectly constant but varies
slightly depending on the intensity (loudness) of the sound wave and the
diameter of the tube it travels through.

\textbf{In Theory of Sound}: Rayleigh cites Regnault's measurements to
bridge the gap between idealized mathematical wave equations and the
messy reality of friction, thermal conduction, and atmospheric pressure
in real-world pipes.

\textbf{The Siren of Baron Charles Cagniard de la Tour} Invented by
French physicist Baron Charles Cagniard de la Tour in 1819, this device
was the first mechanical instrument capable of measuring the exact
absolute frequency (pitch) of a sound.

\textbf{How it worked}: It consisted of a pressurized air chamber with a
fixed metal disk (stator) on top, overlaid by a rotating metal disk
(rotor). Both disks were punched with a series of regularly spaced
holes.

\textbf{The Physics}: Crucially, the holes were drilled at opposing,
slanted angles. When air was forced upward from a bellows, it rushed
through the holes and physically pushed the top disk, making it spin. As
the disk spun, it rapidly cut off and opened the air stream, releasing
regular, rapid \textit{puffs} of air that the human ear perceives as a musical
tone.

\textbf{Rayleigh's Use}: Gears on the spinning disk connected to a dial
mechanism that precisely counted the number of rotations per second. By
multiplying the rotations by the number of holes, scientists could
calculate the exact frequency of the tone mathematically, providing an
invaluable tool for calibrating early acoustics experiments.

\textbf{The Mythological Name}:
Cagniard de la Tour named the device a siren because it could produce
sound underwater.

\textbf{Kaleidophone} Invented by Sir Charles Wheatstone in 1827, the
kaleidophone (literally \textit{beautiful sound-vessel}) was a device designed
to visually demonstrate the paths of vibrating objects using the
phenomenon of persistence of vision.

\textbf{How it worked}: It consisted of a thin, elastic metal rod
clamped firmly at one end. At the free upper end, a small, highly
polished or silvered glass bead was attached to reflect a point of
light.

\textbf{Rayleigh's Context}: When the rod was bowed or struck, it
vibrated transversely in two directions. Because the elasticity of the
rod was usually slightly different in different directions, the bouncing
point of light traced luminous, overlapping, and looping patterns in the
air---known today as Lissajous curves. Rayleigh references it when
discussing the compounding of two perpendicular, independent harmonic
vibrations of varying frequencies.

\textbf{Tonometer} Invented by Johann Heinrich Scheibler in 1834, the tonometer
was a highly precise acoustic apparatus used as an absolute frequency
analyzer and tuning standard before electronic measurement existed.

\textbf{How it worked}: It consisted of a large, sequentially tuned set
of precision tuning forks---often 56 or more---spanning exactly one
octave. Each fork was meticulously filed down so that it was separated
from the previous one by a tiny, identical interval (usually exactly 4
Hz), creating an audible rhythmic beating sound when played together.

\textbf{Rayleigh's Context}: By counting the exact number of beats per
second generated between an unknown acoustic pitch and the closest
tuning forks on the tonometer, scientists could mathematically calculate
the absolute frequency of any sound with extreme accuracy. Rayleigh used
these to study pure tones and verify the exact interference patterns of
wave frequencies.

Here are a few sentences guessing importance of Marin Mersenne
(1588--1648), minimizing almost everything, and otherwise exaggerating: He was
able to get a few hundred of Western Europe's best scientists to adopt
of each others work. He created a clearing house for ideas, centered
around Parisien scientists, and created or standardized Peer Review. In
acoustics, he was among the first to experimentally measure the absolute
speed of sound in air. Mersenne's Laws (1636) mathematically describe
the behavior of vibrating strings. He proved how a string's frequency
changes based on its length, cross-sectional diameter, tension, and
material density. This work serves as Rayleigh's foundation for building
the transverse string vibration equations.

\textbf{Sonometer} Also historically known as a monochord, the sonometer is an
ancient physics apparatus (dating back to Pythagoras) used to study the
mathematical laws of vibrating strings.

\textbf{How it worked}: It features a long, hollow wooden resonance box with one
or more metal wires stretched tightly across it. The wires are anchored
at one end, while the other end uses hanging weights or tuning pegs to
meticulously control the tension. A movable wooden \textit{bridge} can be slid
underneath the string to change its effective vibrating length without
changing its tension.

\textbf{Rayleigh's Context}: Plucking or bowing the string produces a clear tone.
Rayleigh used the sonometer to experimentally prove that the frequency
of a string is inversely proportional to its length, directly
proportional to the square root of its tension, and inversely
proportional to its mass density.

\textbf{Son Rauque}
Savart's \emph{son rauque} F{é}lix Savart's observed \emph{son rauque}
during 1831 experiments on the longitudinal vibrations of solid
cylindrical bars. The main tone of a bar is its longitudinal pitch.
Alongside its main tone the bar unexpectedly emits a rough,
lower-pitched note that is exactly one octave lower. In the 1850s Alfred
Terquem (1831--1887) explained son rauque as an acoustic coupling
effect. The deformation for a High frequency longitudinal
stretch/compression forces a transverse (bending) vibration. This
secondary bending wave acts as a sub-harmonic, vibrating at half the
frequency.

\textit{ Terquem collaborated with instrument makers like Rudolph Koenig}.
Koenig manufactured Terquem's \emph{verges de laiton} adopted by
the music education curriculum.

P.S.  This might be an AI hallucination.
I found no references to inventions of Terquem.
Alfred Terquem was an experimental physicists.
Koenig learning from Helmholtz to make H's apparatus',
such as resonators.
Then Koenig fabricated the sound flame, which was crucial.
I did not look for Savart's son rauque or \emph{verges de laiton}.


\textbf{J. Bourget} (1822--1887) fluids and solids, circular membrane vibration (1866),

section to the vibrations of membranes. He explicitly quotes the
collaborative experimental work of Bourget and Bernard, who tested how
circular membranes react when physically blasted with acoustic notes
from organ pipes.

Chladni's plates or diagrams? Mathieu?

Rudolph Koenig (1832--1901) transformed acoustics from a theoretical branch of
math into a precise experimental science.  He invented the manometric flame
apparatus which visualized sound waves using gas flames.  Also, he perfected
the construction of ultra-precise tuning forks and tonometers.  Rayleigh relied
on a manometric flame apparatus in his experiments.  Rayleigh cites Koenig's
empirical data on interferents between subtle overtones, and changes of pitch
in tuning forks near a resonant cavity,

\section{Torsional Vibration}\label{torsional-vibration}

When a uniform, circular bar or shaft experiences twisting along its
longitudinal axis, it undergoes torsional vibration.


\begin{longtable}[]{@{}cc@{}}
\tabularnewline%
quantity & symbol \\
\midrule\noalign{}
\endfirsthead%
\toprule\noalign{}
quantity & symbol \\
\midrule\noalign{}
\endhead%
\endlastfoot{}%
axial position & \(x\)  \\
 angle of twist & \(\theta (x,t)\) \\
torsional wave velocity & \(c = \sqrt{\frac{G}{\rho}}\) \\
\end{longtable}

Torsional vibrations share an the mathematical structure with
axial/longitudinal vibrations, with \(\theta\) instead of \(u\).
Assuming a circular cross-section, e.g., cross-sections remain perfectly
planar, the governing partial differential equation is a standard
one-dimensional wave equation:
\[
\ddot{\theta}=\frac{G}{\rho }\theta_{,xx}.
\]
In the same bar, the ratio between longitudinal and torsional wave speeds
\(\sqrt{2(1+\mu)}\).

\section{Thermoacoustics}\label{thermoacoustics}

Pieter Rijke discovered the Rijke Tube (Rye-kuh, 1859).
Singing flames and the Rijke tube phenomenon
are two variations of thermoacoustics (turning heat into sound).
Both rely on a physics principle known as the Rayleigh Criterion, which
states that if heat is added to a gas at the moment of highest
compression (maximum pressure), the acoustic vibration will be
amplified.

The Singing Flame (Higgins' Flame, 1777)

Setup: A small gas flame (originally hydrogen) is placed inside a larger
vertical glass or metal tube open at both ends.

The Mechanism: The heat from the flame causes air to expand and rush up
the tube. This airflow naturally triggers a weak, native acoustic
standing wave inside the pipe. The pressure changes of this wave force
the flame to flicker violently. Because the flame fluctuates in perfect
synchronization with the air pressure, it releases bursts of heat at
exactly the right moments to amplify the sound wave, causing the pipe to
loudly \textit{sing}.

The Rijke Tube (1859), Dutch physicist Pieter Rijke

Setup: Instead of an open flame inside the pipe, a fine metal wire mesh
is placed inside the lower half of a vertical tube.

The Mechanism: The mesh is heated with a burner until it is red-hot. The
burner is then removed.

The Result: Convection draws cold air into the bottom of the tube. As
this air passes up through the hot mesh, the mesh transfers heat to the
air in periodic pulses that sync with the pipe's natural resonant
frequency, creating an incredibly loud, deep, foghorn-like roar without
any active flame present.

In short: A singing flame uses an active, oscillating chemical reaction
to generate thermoacoustic sound, whereas a Rijke tube uses a static,
pre-heated metal element combined with buoyant convection to achieve the
exact same acoustic amplification.

\section{Appendix}\label{appendix}

The discovery of sub-harmonics in fluids, Faraday waves, motivated the
discovery of acoustic sub-harmonics. Amazing discussion of Faraday
waves, and related period states.

Professor Ratul Dasgupta, Introduction to Interfacial Waves, Lectures
\( 1,2, \dots 52, \ldots \)

\url{https://www.youtube.com/watch?v=VJ-Me2Q3tdg}

\url{https://www.youtube.com/watch?v=J6i4izpfEtw}

\url{https://www.youtube.com/watch?v=2BOGnDg9NX8}

dead link: \url{https://nptel.ac.in/course.html}

More vocabulary.

\begin{itemize}
\item
  aliquot: an `aliquot part' of an integer is an exact divisor or
  submultiple.
\item
  gutta-percha: thermoplastic latex sap.
\item
  gutta-percha tube: a hollow pipe used for communication by sound
  (stethoscopes until 1950)
\item
  diapson: sound made by a pipe organ when a chord is held.
\item
  Faulhorn: a mountain in the Alps
\item cateris paribus: all other things being equal
\item concord: a harmonious, stable combination of notes that sounds complete
and at rest to the ear.
\item ex hypothesi: according to the hypothesis proposed
\item Leyden : capacitor
\end{itemize}

Article 19. The sound is like the howl of a wolf, and is called the wolf
interval.

Capacitance is the shunt capacitance per unit length.
Conductance represents the imperfect insulation between wires, where a tiny
amount of current \textit{leaks} through the dielectric material.

contemporary RGLC parameters,  \( j = \sqrt{-1} \)

\begin{tabular}{cc}
\(\Omega \) & Ohm \\
\(H\) &  Henry  \\
\(F\) &  Farad  \\
\(R\) & \(\Omega/m\)  \\
\(G\) &  Leakage-Conductance \\
 &  Conductance (S/m) \\
\(L\) &  inductance \(H/m\) \\
\(C\) & permittance \(F/m\) \\
 & capacitance  \\
\(S\) & Siemen \\
\end{tabular}

I do not know what notation Rayleigh used.

Voltage Drop and Current Leakage equations,
\[
\frac{\partial}{\partial x}
\begin{bmatrix} V \\ I \end{bmatrix}
+
\frac{\partial}{\partial t}
\begin{bmatrix}
0 & L \\
C & 0
\end{bmatrix}
\begin{bmatrix} V \\ I \end{bmatrix} = -
\begin{bmatrix}
0 & R \\
V & 0
\end{bmatrix}
\begin{bmatrix} V \\ I \end{bmatrix}
\]

Impedance
\[
Z_{0}(\omega) =\sqrt{\frac{R+j\omega L}{G+j\omega C}}
\]

The L in LCR meter refers to an inductometer.

Vocabulary:
hundred weight (cwt), 112 lbs (force) =  1 hundred weight

\[\vec{x} = r( t - \sin t, 1 - \cos t)^T \tag{cycloid} \]
\[\vec{x} = ( r t - d \sin t, r - d \cos t)^T \tag{trochoid} \]
curtate \(d < r \) , or prolate \(d > r \) .

aeolian harp, determinental, determinate, extensible, isochronism (constant period),
Germansilver (Schlage keys), crispation, perigee, lycopodium powder (Chladni)

Smee cell  - a type of battery  circa 1900


%%\chapter{Vibrations of Plates}\label{vibrations-of-plates}

214.\footnote{This Chapter deals only with flexural vibrations. The
extensional vibrations of an infinite plane plate are briefly considered in
Chapter X.a, as a particular case of those of an infinite cylindrical shell.
They are not of much acoustical importance.}
In order to form according to Green's method the equations of
equilibrium and motion for a thin solid plate of uniform isotropic
material and constant thickness, we require the expression for the
potential energy of bending. It is easy to see that for each unit of
area the potential energy V is a positive homogeneous symmetrical
quadratic function of the two principal curvatures, Thus, if p,, p, be
the principal radii of curvature, the expression for V will be

where A and yw are constants, of which A must be positive, and must be
numerically less than unity. Moreover if the material be of such a
character that it undergoes no lateral contraction when a bar is pulled
out, the constant must vanish. This amount of information is almost all
that is required for our purpose, and we may therefore content ourselves
with a mere statement of the relations of the constants in (1) with
those by means of which the elastic properties of bodies are naturally
defined.

\todo[inline]{The smudged numerator of Thomson's experssion \(E(\kappa,n)\)
is denoted here by \(\Upsilon\)}

From Thomson and Tait's Natural Philosophy, \S 639, 642, 720, it
appears that, if b be the thickness, g Young's modulus, and yw the ratio
of lateral contraction to longitudinal elongation when a bar is pulled
out, the expression for V is
\footnote{The following comparison of the notations for Young's modulus and the ratio of lateral contraction to longitudinal elongation used by the principal writers may save trouble to those who wish to consult the original memoirs.
\begin{tabular}{cc}
Clebsch &\(E\) \\
Thomson &\(M\) \\
Thomson &\(\frac{\Upsilon }{3 \kappa +n }\) \\
Thomson &\(\frac{n(3m-n)}{m}\) \\
Kirchhoff,Donkin &\(q\) \\
Kirchhoff &\(2K \frac{1+3 \theta}{1+ \theta}\) \\
\end{tabular}
\begin{tabular}{cc}
Clebsch,Donkin &\(\mu\) \\
Thomson &\(\sigma\) \\
Thomson &\(\frac{m-n}{2m}\) \\
Kirchhoff &\(\frac{\theta}{1+2 \theta}\) \\
\end{tabular}
}.

[Equation (2) gives the interpretation of the constants of (1)
in its application to a homogeneous plate of isotropic material ;
but the expression (1) itself is of fai- wider scope. The material
composing the plate may vary from layer to layer, and the elastic
character of any layer need not be isotropic, but only symmetrical
with respect to the normal. As a particular case, the middle
layer, or indeed any other layer, may be supposed to be physically
inextensible.

Similar remarks apply to the investigations of the following
chapter relating to curved shells.]

If w be the small displacement perpendicular to the plane of the plate
at the point whose rectangular coordinates in the plane of the plate are
x, y,

and thus for a unit of area, we have

V= oy ah os viv)" - 2 (1-H) oa (zer,) (3),

which quantity has to be integrated over the surface (8) of the plate.

215. We proceed to find the variation of \(V\), but it should be previously
noticed that the second term in \(V\), namely
\[
\iint \frac{ dS }{ \rho_1 \rho_2 }
\]
represents the total curvature of the plate, and is therefore dependent
only on the state of things at the edge.
\[ \delta V = \frac{1}{1} \iint dS;\]
(1) so that we have to consider the two variations
\[\iint \Lap w \Lap \delta w dS,\quad \iint \delta \frac{ dS }{ \rho_1 \rho_2 }.\]

Now by Green's theorem
\[\iint \Lap w \Lap \delta w dS=
\]

(2) in which ds denotes an element of the boundary, and g denotes
differentiation with respect to the normal of the boundary drawn
outwards.

The transformation of the second part is more difficult. We have
\[\iint \delta \frac{ dS }{ \rho_1 \rho_2 }=\]

The quantity under the sign of integration may be put into the form

\(\forall\) smooth \(F(x,y)\),
\begin{equation}
\iint_D \left(\frac{\del}{\del x},\frac{\del}{\del y}\right)F \, dx \, dy=
\int_{\del D} F (\sin(theta), \cos(theta)) ds
\end{equation}

(3) \(\theta = \angle \) between x and the normal drawn outwards, and the
integration on the right-hand side extends round the boundary. Using
these, we find

If we substitute for ' dy their values in terms Tn Buy, from the
equations (see Fig. 40) (4)

Fig. 40.
we obtain

ddw i 2a'w dw 3 5 - = as a 9 FY +. (OG a

+ fa om 8 in o (3 - Gp) i (rt). (5)

The second integral by a partial integration with respect to s may be
put into the form
\[ Jou J, eos O (0 (32 - Siz) + ( \cot() - in 2 6) - \]
Collecting and rearranging our results, we find 9 80= isd dy w tw dw -
weds F +(1- 1) J (008 8 (Se) - 53) . 29 in a g, ew + (e) (ay) + ds eve +
(1-4)) ( (to 5 + (e) +2 () (rs) -. (6). There will now be no difficulty
in forming the equations of
motion. If p be the volume density, and Z b dS the transverse force
acting on the element dS,

V - Z b w dS + b w w dS =0 (7)

is the general variational equation, which must be true whatever
function (consistent with the constitution of the system) dw may be
supposed to be. Note that the rotary inertia is neglected in equation
(7). Hence by the principles of the Calculus of Variations
at every point of the plate. If the edges of the plate be free, there is
no restriction on the
hypothetical boundary values of dw and Bu , and therefore the
coefficients of these quantities in the expression for 6 V must vanish.

The conditions to be satisfied at a free edge are thus
\[a W P dw (Vey - 1) G, 008 \sin(Fa) +  \cos(oe) Pw  =0 : Pw Ge dy da.\]
(9).
\[(1-p) + \sin(a)  \sin(Ai);\]
If the whole circumference of the plate be clamped, Sw = 0, aw =0,
and the satisfaction of the boundary conditions is already scoured. If
the edge be `supported'\footnote{Compare Article 162}, w =0, but SY is
arbitrary. The second of the equations (9) must in this case be satisfied by w.

216. Conditions for a free edge.
The boundary equations may be simplified by getting rid of the
extrinsic element involved in the use of Cartesian coordinates. Taking
the axis of parallel to the normal of the bounding curve, we see that we
may write

(t) (agp) tee in 8 = ou aw a Also OW TS cette (1),

where o is a fixed axis coinciding with the tangent at the point under
consideration. In general differs from ae To obtain the relation between
them, we may proceed thus. Expand w by Maclaurin's theorem in ascending
powers of the small quantities n and o, and substitute for x and o their
values in terms of s, the arc of the curve.

Thus in general

dw dw dw dw, 7

w= w+ on, n+ ote + and a," +4 age +

while on the curve =8 + cubes, mended where p is the
radius of curvature. Accordingly for points on the curve,
and therefore do o1dw ; ds do Pp dn Ce reverses (3) whence from (1)

We conclude that the second boundary condition in (9) Article 215 may be
put into the form

In the same way by putting = 0, we see that

\[\cos(\theta) \sin(G) + ( \cos(6 - \sin(oe)\]

is equivalent to ou where it is to be understood that the axes
of n and o are fixed. The first boundary condition now becomes
d d w = Vw +(1- 1) 5(a) = 0 (5).

If we apply these equations to the rectangle whose sides are parallel to
the coordinate axes, we obtain as the conditions to be satisfied along
the edges parallel to y,

d/dat (dw) dw/ays ae art O- ap (6).

In this case the distinction between o and s disappears, and p, the
radius of curvature, is infinitely great. The conditions for the other
pair of edges are found by interchanging x and y. These results may be
obtained equally well from (9) Article 215 directly, without the
preliminary transformation.

217. Conjugate Property.  If we suppose Z =0, and write

the general equation becomes

or, if (pt-e), (3), where (4).

Any two values of w, w and v, corresponding to the same boundary
conditions, are conjugate, that is to say .

u v dS =0 (5),

provided that the periods be different. In order to prove this from the
ordinary differential equation (3), we should have to retrace the steps
by which (8) was obtained. This is the method adopted by Kirchhoff for
the circular disc, but it is much simpler and more direct to use the
variational equation

V + b w w dS =0 (6),

in which w refers to the actual motion, and w to an arbitrary
displacement consistent with the nature of the system. V is a
symmetrical function of w and dw, as may be seen from Article 215, or
from the general character of V (Article 94.)

If we now suppose in the first place that w=u, w=v, we have

V = pop' u v dS;

and in like manner if we put w =v, dw =u, which we are equally entitled
to do,

8V = pbp'' , whence

(p -p'') =0 (7).

This demonstration is valid whatever may be the form of the boundary,
and whether the edge be clamped, supported, or free, in whole or in
part.

As for the case of membranes in the last Chapter, equation (7) may be
employed to prove that the admissible values of p'' are real; but this
is evident from physical considerations.

218. Polar Coordinates.
For the application to a circular disc, it is necessary to express
the equations by means of polar coordinates. Taking the centre of the
disc as pole, we have for the general equation to be satisfied at all
points of the area
\( -\kappa 4) w=0\)
(1), where (Article 200)

To express the boundary condition (Article 216) for a free edge (r=a) we
have

= radius of curvature = ; and thus

(2).

After the differentiations are performed, r is to be made equal to a.

If w be expanded in Fourier's series W=WU+U, +... +, +...,

each term separately must satisfy (2), and thus, since

w,x (n) - a, - (Mw, ldw. /2 -pdw, 3-h ie (ae oa) nt (a dr a w,)=0

The superficial differential equation may be written (y+) (v-) w=0,
which becomes for the general term of the Fourier expansion

shewing that the complete value of w, will be obtained by adding
together, with arbitrary constants prefixed, the general solutions of

(4).

The equation with the upper sign is the same as that which obtains in
the case of the vibrations of circular membranes, and as in the last
Chapter we conclude that the solution applicable to the problem in hand
is w, cc J, (r), the second function of r being here inadmissible. In
the same way the solution of the equation with the lower sign is w, J,
(ier), where = / -1 as usual.
[See § 221 a.]
The simple vibration is thus

\[w, = \cos(nd a J, (er) + Bd, (ter) + \sin(n) J, (xr) + SJ, (ter)\]

The two boundary equations will determine the admissible values of and
the values which must be given to the ratios a: and y: From the form of
these equations it is evident that we must have a:B=7:6,

and thus w, may be expressed in the form

\[w, = P \cos(nb - a) J, (er) + Xd, (ter) \cos(pt-e)\]

(5).

As in the case of a membrane the nodal system is composed of the n
diameters symmetrically distributed round the centre, but otherwise
arbitrary, denoted by

((MO - 2) =0 ...(6), together with the concentric circles, whose
equation is Ti (ar) +N, (tH) =0 (7).

219. Kirchhoff's Theory. % 308
In order to determine and we must introduce the boundary
conditions, When the edge is free, we obtain from (3) Article 218,

in which use has been made of the differential equations satisfied by J,
(r), J, (ier). In each of the fractions on the right the denominator may
be derived from the numerator by writing e in place of x. By elimination
of the equation is obtained whose roots give the admissible values of x.
,

When x = 0, the result assumes a simple form, viz.

2 (1-p) + xa ye + 0a 3 0 (2). This, of course, could have been more
easily obtained by neglecting n from the beginning.

The calculation of the lowest root for each value of n is troublesome,
and in the absence of appropriate tables must be effected by means of
the ascending series for the functions J, (r), J, (tr). In the case of
the higher roots recourse may be had to the semiconvergent descending
series for the same functions, Kirchhoff finds

\[ Blo Dp, \tan((xa - nm) = S8 Sea oa) \]
(3), where A=y=(1-4) B =y (1- 4n') -8, C =x (1 - 4n) (9 - 4n ) + 48 (1 +
4n'), = -'y (1 - 4n ) (9 - 4n ) (18 - 4n) + 8 (9 + 186n + 80x).


When a is great,

tan (ca-4 n r) =0 approx. ; whence (4), where h is an integer.

It appears by a numerical comparison that A is identical with the number
of circular nodes, and (4) expresses a law discovered by Chladni, that
the frequencies corresponding to figures with a given number of nodal
diameters are, with the exception of the lowest, approximately
proportional to the squares of consecutive even or uneven numbers,
according as the number of the diameters is itself even or odd. Within
the limits of application of (4), we see also that the pitch is
approximately unaltered, when any number is subtracted from A, provided
twice that number be added ton. This law, of which traces appear in the
following table, may be expressed by saying that towards raising the
pitch nodal circles have twice the effect of nodal diameters. It is
probable, however, that, strictly speaking, no two normal components
have exactly the same pitch.

The table, extracted from Kirchhoff's memoir, gives the pitch of the
more important overtones of a free circular plate, the gravest being
assumed to be C' . The three columns under the heads Ch, P, W refer
respectively to the results as observed by Chladni and as calculated
from theory with Poisson's and Wertheim's values of . A plus sign
denotes that the actual pitch is a little higher, and a minus sign that
it is a little lower, than that written.

Dangling\footnote{ Gis corresponds to G of the English notation, and h to b natural. .}

The discrepancies between theory and observation are considerable, but
perhaps not greater than may be attributed to irregularity in the plate.

220.  Oscillation of Nodes.
The radii of the nodal circles in the symmetrical case (n=0) were
calculated by Poisson, and compared by him with results obtained
experimentally by Savart. The following numbers are taken from a paper
by Strehlke', who made some careful measurements. The radius of the disc
is taken as unity.

Observation. Calculation. One circle ... 067815 068062.

039133 039151.

Two circles... e149 0'84200.

025631 025679. Three circles 059107 059147. 089360 089381.

The calculated results appear to refer to Poisson's value of u, but
would vary very little if Wertheim's value were substituted.

The following table gives a comparison of Kirchhoff's theory (n not zero)
with measurements by Strehlke made on less accurate discs.

Radius of Circular Nodes.

Calculation. P=7(P). 2=s (W.). 0-781 0783 0-781 0-783 0-78136 0-78088

Observation.

0:79 081 082 082194 082274 0838 0842 0-84523 084681 0-488 0-492 0-49774
049715 0869 0-869 0-87057 087015

Dangling\footnote{ Pogg. Ann. xcv, p. 577, 1855.}

221. Oscillation of Nodes. % 365
When the plate is truly symmetrical, whether uniform or not, theory
indicates, and experiment verifies, that the position of the nodal
diameters is arbitrary, or rather dependent only on the manner in which
the plate is supported. By varying the place of support, any desired
diameter may be made nodal. It is generally otherwise when there is any
sensible departure from exact symmetry. The two modes of vibration,
which originally,
in consequence of the equality of periods could be combined in any
proportion without ceasing to be simple harmonic, are now separated and
affected with different periods. At the same time the position of the
nodal diameters becomes determinate, or rather limited to two
alternatives. The one set is derived from the other by rotation through
half the angle included between two adjacent diameters of the same set.
This supposes that the deviation from uniformity is small; otherwise the
nodal system will no longer be composed of approximate circles and
diameters at all. The cause of the deviation may be an irregularity
either in the material or in the thickness or in the form of the
boundary. The effect of a small load at any point may be investigated as
in the parallel problem of the membrane Article 208. If the place at
which the load is attached does not lie on a nodal circle, the normal
types are made determinate. The diametral system corresponding to one of
the types passes through the place in question, and for this type the
period is unaltered. The period of the other type is increased.


[The divergence of free periods, which is due to slight inequalities, would
seem to afford an explanation of some curious observations by Savart When a
circular plate, vibrating with nodal diameters, is under the influence of the
bow applied at any part of the circumference, the nodal diameters indicated by
an arc so situated that the bow lies in the middle of a vibrating segment.
If, however, the bow be suddenly withdrawn, the nodal system
oscillates, or even revolves, during the subsidence of the motion.
It is evident that no such displacement could be expected,
were the plate absolutely symmetrical. The same would be true,
even in the case of asymmetry, if the bow were so applied as to
excite one only of the two determinate vibrations then possible.
But in general the effect of the bow must be to excite both kinds
of vibrations, and then the matter is more complicated. It would
seem that so long as the constraining action of the bow lasts, both
vibrations are forced to keep the same time, and the effect is
much the same as in the case of symmetry. But on withdrawal
of the bow the free vibrations which then ensue take place each in
its proper frequency, and a phase difference soon arises by which
the effects are modified.

Let us suppose that the origin of 0 is so chosen in relation to the
irregularities that the types of vibration are represented by
cos nO, sin n0. Then in general the free vibrations, resulting from the action
of the bow at an arbitrary point of the circumference, may be taken to be
\[
\cos( na ) \sin (\beta)  \cos \pi - \sin \beta ia \cos(  \beta t)  \cos (pf -h e)
\]
(1), where e is the difference of phase which is accumulated since
the commencement of the free vibrations. In the case
symmetry e = 0, and (1) becomes

sin n(0 - a) cos pt (2),
Dangling\footnote{Ann. Chim., vol. 86, p. 257, 1827.}

which represents a fixed nodal system

 = a + r i(v n) (3),
in any arbitrary position depending upon the point of application
of the bow. A similar fixity of the nodal system occurs, in spite
of the variable e, when a is so chosen that \( \cos noL = 0\) or  \( \sin wa = 0\).
But in general there is no fixed nodal system. When e is a
multiple of 27r, that is when the two vibrations are restored to
the same phase, there is a nodal system represented by (3). And
when € is an odd multiple of tt, so that the two vibrations are in
opposite phases, we have in place of (2)
sin n(0 -h a) co pt (4),

with a nodal system

= - a 4- m (tt w) (.5).
In these cases there is a nodal system, and in a sense the system
may be said to oscillate between the positions given by (3) and (5) ;
but it must not be overlooked that at intermediate times there is
no true nodal system at all. Thus, when e = Jtt, (1) becomes

cos no sin n cos pt + sin wa cos nd sin pt

The squared amplitude of this motion is

cos nd sin nO H- sin na cos nO,

a quantity which does not vanish for any value of 6, In general
the squared amplitude is

cos na sin rtO + sin na cos 7)0 - 2 cos na sin na cos n0 sin n0 cos e,
or, as it may also be written,
\[ \cos 2na \cos 2n0 -  \sin 2wa \sin 2n0 \cos e.\]

(6).  This quantity is a maximum or a minimum when

tan 2,n0 = cos € tan 2na (7).
The minimum of motion thus oscillates backwards and forwards
between = +a and = -a; but as we have seen, it is only in
these extreme positions that the minimum is zero.

A like phenomenon occurs during the free vibrations of a
circular membrane, or in fact of any system of revolution such
that the position of nodal lines is arbitrary so long as the
symmetry is complete.]

The most general motion of the uniform circular plate is expressed by
the superposition, with arbitrary amplitudes and phases, of the normal
components already investigated. The determination of the amplitude and
phase to correspond to arbitrary initial displacements and velocities is
effected precisely as in the corresponding problem for the membrane by
the aid of the characteristic property of the normal functions proved in
Article 217.

The two other cases of a circular plate in which the edge is either
clamped or supported would be easier than the preceding in their
theoretical treatment, but are of less practical interest on account of
the difficulty of experimentally realising the conditions assumed. The
general result that the nodal system is composed of concentric circles,
and diameters symmetrically distributed, is applicable to all the three
cases.

\section{Clamped boundary}\label{sect:221a}
221 a.
The use in the telephone of a thin circular plate clamped at the
edge lends a certain interest to the calculation of the periods and
modes of vibration of such a plate. It will suffice to consider the
symmetrical modes.

By (5) § 218 we may take as representing the in this case
\[ w = J_o (kr) + \lambda J_o(ikr)  = Jo (kr) + \lambda I_o(kr) \]
from which
\[ \frac{1}{k} \frac{d}{dr} w = -J_1 (kr) + \lambda I_1(kr)\]
where we write
\[ I_o(z) = J_o(iz), I_1(z) = i J_oT{\prime}(iz).\]
Since the plate is clamped at \(r = a\), both \(w\) and \(dw/dr\) must
there vanish. Hence, writing \(ka = z\), we get as the frequency equation
\[ \frac{J_1(z)}{J_o(z)} + \frac{I_1(z)}{I_o(z)} = 0.\]
.(5). In (5) \(I_1\) and \(I_o\) are both positive,
so that the signs of \(J_i\) and \(J_o\)
must be opposite. Hence by Table B § 206 the first root must lie
between \(2.4\) and \(3.8\), the second between \(5.5\) and \(7.0\), and
so on. The values of the earlier roots might be obtained without
much difficulty from the series for \(I_o\) and \(I_1\)
by using the table § 200 for \(J_o\) and \(J_1\); but
it will be convenient for the present and further purposes to give a
short table \footnote{Calculated by A. Lodge, Brit. Ass. Rep., 1889.}
of the functions \(I_o\) and \(I_1\)
themselves. For large values of the argument descending series,
analogous to (10) § 200, may be employed.

\todo[inline]{ \(I_o\) and \(I_1\) are tabulated on \( [0,6] \). }
% page 367 top

The first root of (5) is \(z=3.20\). This then is the value of \(ka\) for the
gravest symmetrical vibration. The next value of z is about 6.3. Since
the frequency varies as \(k^2\) (§ 217), the interval
between the tones is nearly two octaves.

Returning to the first root, we have for the frequency \(n\) §217,
\[ n = p/(2 \pi) \]
This is the general formula. For rough calculations \(\mu^2\) in the
denominator may be omitted. If for the case of iron we take
\[ \rho = 7.7, q = 2.0 \times 10^{12}\]
we find \(n  = 2.4 \times 10^{5}  2h/a^2\) (7),
2h and a being expressed in centimetres.

A telephone plate measured by the author gave
\[ a = 2.2, 2h = 0.020.\]
According to these values \(n = 991\) vibrations per second.
222. We have seen that in general Chladni's figures as traced by sand
agree very closely with the circles and diameters of theory ; but in
certain cases deviations occur, which are usually attributed to
irregularities in the plate. It must however be remembered that the
vibrations excited by a bow are not strictly speaking free, and that
their periods are therefore liable to a certain modification. It may be
that under the action of the bow two or more normal component vibrations
coexist. The whole motion may be simple harmonic in virtue of the
external force, although the natural periods would be a little
different. Such an explanation is suggested by the regular character of
the figures obtained in certain cases.

Another cause of deviation may perhaps be found in the manner in which
the plates are supported. The requirements of theory are often difficult
to meet in actual experiment. When this is 80, we may have to be content
with an imperfect comparison; but we must remember that a discrepancy
may be the fault of the experiment as well as of the theory.

[In the ordinary use of sand to investigate the vibrations of
flat plates and membranes the movement to the nodes is irregular
in its character. If a grain be situated elsewhere than at a node,
it is made to jump by a sufficiently vigorous transverse vibration.
The result may be a movement either towards or from a node;
but after a succession of such jumps the grain ultimately finds its
way to a node as the only place where it can remain undisturbed.
Grains which have already arrived at a node remain there, while others are
constantly shifting their position.

It was found by Savart that very fine powder, such as lycopodium, behaves
differently from sand. Instead of collecting at
the nodes, it heaps itself up at the places of greatest motion.
This effect was traced by Faraday to the influence of currents of
air, themselves the result of the vibration. In a vacuum all
powders move to the nodes.

In some cases the movement of sand to the nodes, or to some
of them, takes place in a more direct manner as the result of
friction. Thus, in his investigation of the longitudinal vibrations
of thin narrow strips of glass, held horizontally, Savart observed
the delineation of nodes apparently dependent upon an accompaniment of
vibrations of a transverse character. The special
peculiarity of this phenomenon was the non-correspondence of the
lines traced by sand upon the two faces of the glass when tested
\footnote{On a Familiar Class of Acoustical Figures, Phil Trans. 1831, p. 299.
Ann. d. Chim.t vol. 14, p. 118, 1820.}


223. The first attempt to solve the problem with which we have just been
occupied is due to Sophie Germain, who succeeded in obtaining the
correct differential equation, but was led to erroneous boundary
conditions. For a free plate the latter part of the problem, is indeed
of considerable difficulty. In Poisson's \emph{Mémoire sur l'équilibre}
\emph{et le mouvement des corps élastiques}, that eminent mathematician
gave three equations as necessary to be satisfied at alt points of a
free edge, but Kirchhoff has proved that in general it would be
impossible to satisfy them all. It happens, however, that an exception
occurs in the ease of the symmetrical vibrations of a circular plate,
when one of the equations is true identically, Owing to this
peculiarity, Poisson's theory of the symmetrical vibrations is correct,
notwithstanding the error-in his view as to the boundary conditions.
In 1850 the subject was resumed, by Kirchhoff', who first gave the two
equations appropriate to 9 free edge, and completed the theory of the
vibrations of a circular disc.

224. The correctness of Kirchhoff's boundary equations has been disputed
by Mathieu'; who, without explaining where he considers Kirchhoff's
error to He, has substituted a different set of equations. He proves
that if w and u' be two normal functions, so that w =u (pt), w=wu' (p't
are possible vibrations, then

(1)

Dangling\footnote{ Mém. de l' Acad. d. Sc. a Par. 1829.}

2 Crelle, t. XL. p. 51, Ueber das Gleichgewicht und die Bewegung einer
elastichen Scheibe,

3 Liouville, t. XIV. 1869,

This follows, if it be admitted that u, u' satisfy respectively the
equations

\[c \Grad u = p u, c \Grad u' = p' u'.\]

Since the left-hand member is zero, the same must be true of the
right-hand member; and this, according to Mathieu, cannot be the case,
unless at all points of the boundary both wu and wu' satisfy one of the
four following pairs of equations:

=0 yu =0 tat .=0 = du ' dy u , o al dy u

The second pair would seem the most likely fora free edge, but it is
found to lead te an impossibility, Since the first and third pairs are
obviously inadmissible, Mathieu concludes that the fourth pair of
equations must be those which really express the condition of a free
edge. In his belief in this result he is not shaken by the fact that the
corresponding conditions for the free end of a bar would be

du 0, a =0,

the first of which is contradicted by the roughest observation of the
vibration of a large tuning fork.

The fact is that although any of the four pairs of equations would
secure the evanescence of the boundary integral in (1), it does not
follow conversely that the integral can be made to vanish in no other
way; and such a conclusion is negatived by Kirchhoff's investigation.
There are besides innumerable other cases in which the integral in
question would vanish, all that is really necessary being that the
boundary appliances should be either at rest, or devoid of inertia.

225. The vibrations of a rectangular plate, whose edge is
supported, may be easily investigated theoretically, the normal
functions being identical with those applicable to' a membrane of the
same shape, whose boundary is fixed. If we assume

w = () = () ((le). (1),
we see that at all points of the boundary, dw dw a = 0, dy = 0,
which secure the fulfilment of the necessary conditions (Article 215).
The value of p, found by substitution in c yw =p',
is p=Cn (5 + ) (2),
shewing that the analogy to the membrane does not extend to the sequence
of tones.

w=0,

It is not necessary to repeat here the discussion of the primary and
derived nodal systems given in Chapter IX. It is enough to observe that
if two of the fundamental modes (1) have the same period in the case of
the membrane, they must also have the same period in the case of the
plate. The derived nodal systems are accordingly identical in the two
cases.

The generality of the value of w obtained by compounding with arbitrary
amplitudes and phases all possible particular solutions of the form (1)
requires no fresh discussion.

Unless the contrary assertion had been made, it would have seemed
unnecessary to say that the nodes of a supported plate have nothing to
do with the ordinary Chladni's figures, which belong to a plate whose
edges are free.

The realization of the conditions for a supported edge is scarcely
attainable in practice. Appliances are required capable of holding the
boundary of the plate at rest, and of such a nature that they give rise
to no couples about tangential axes. We may conceive the plate to be
held in its place by friction against the walls of a cylinder
circumscribed closely round it.


226. Rectangular Plate. page 373
The problem of a rectangular plate, whose edges are free, is one of
great difficulty, and has for the most part resisted attack. If we
suppose that the displacement w is independent of y, the general
differential equation becomes identical with that with which we were
concerned in Chapter VI. If we take the solution corresponding to the
case of a bar whose ends are free, and therefore satisfying
dw dw
T= 0
when x=0 and when 2=a, we obtain a value of w which satisfies the
general differential equation, as well as the pair of boundary equations

me a oe dat +h Gye =

which are applicable to the edges parallel to y; but the second boundary
condition for the other pair of edges, namely

apt eae O (Q) will be violated, unless zw=0. This shews that, except in
the ease reserved, it is not possible for a free rectangular plate to
vibrate after the manner of a bar; unless indeed as an approximation,
when the length parallel to one pair of edges is so great that the
conditions to be satisfied at the Second pair of edges may be left out
of account,

[The case where two opposite edges are free while the other two edges are
supported, been discussed by Voigt (Gottingen Nachrichient ls93) p. I

Although the constant (which expresses the ratio of lateral contraction
to longitudinal extension when a bar is drawn out) is positive for every
known substance, in the case of a few substances - cork, for example - it
is comparatively very small. There is, so far as we know, nothing absurd
in the idea of a substance for which w vanishes. The investigation of
the problem under this condition is therefore not devoid of interest,
though the results will not be strictly applicable to ordinary glass or
metal plates, for which the value of w is about 3.7

If u,, u,, etc. denote the normal functions for a free bar investigated
in "Chapter VIII, corresponding to 2, 3, nodes, the vibrations of
a rectangular plate will be expressed by

w=u, (3), w=u,(2), etc., or w= u(.), w= u, (8), etc.

Dangling\footnote{ In order to make a plate of material, for which is not zero, vibrate
in the manner of a bar, it would be necessary to apply constraining
couples to the edges parallel to the plane of bending to prevent the
assumption of a contrary curvature. The effect of these couples would be
to raise the pitch, and therefore the calculation founded on the type
proper to 1=0 would give a result somewhat higher in pitch than the
truth}

In each of these primitive modes the nodal system is composed of
straight lines parallel to one or other of the edges of the rectangle.
When b=a, the rectangle becomes a square, and the vibrations w=, (2),
w=u, (2),

having necessarily the same period, may be combined in any proportion,
while the whole motion still remains simple harmonic, Whatever the
proportion may be, the resulting nodal curve will of necessity pass
through the points determined by

4, (2) =0, 4, ( ) =0.

Now let us consider more particularly the case of n=1. The nodal system
of the primitive mode, w=u, (2), consists

of a pair of straight lines parallel to y, whose distance from the
nearest edge is 2242a. The points in which these lines are met - by the
corresponding pair for w =u, (4), are those through which the nodal
curve of the compound vibration must in all cases pass, It is evident
that they are symmetrically disposed on the diagonals of the square. If
the two primitive vibrations be taken equal, but in opposite phases
(or, algebraically, with equal and opposite amplitudes), we have ,

w= Ys, ()

from which it is evident that w vanishes when a=y, that is along the
diagonal which passes through the origin. That w will also vanish along
the other diagonal follows from the symmetry of the functions, and we
conclude that the nodal system of (3) Fig. 41.


comprises both the diagonals (Fig. 41). This is a well-known mode of
vibration of a square plate.

A second notable case is when the amplitudes are equal and their phases
the same, so that w=U, (3) +4, () (4)

The most convenient method of constructing graphically the curves, for
which \(w\) is constant, is that employed by Maxwell in similar cases. The two
systems of curves (in this instance straight lines) represented by wu,
\(G\) is constant, first laid down, the values of the constants forming an
arithmetical progression with the same common difference in the two
cases. In this way a network is obtained which the required curves cross
diagonally. The execution of the proposed plan requires an inversion
of the table given in Chapter VL, Article 178, expressing the march of
the function u, of which the result is as follows :

The system of lines represented by the above values of x
(completed symmetrically on the further side of the central line) and
the corresponding system for y are laid down in the figure (42). From
these the curves of equal displacement are deduced. At the centre of the
square we have w maximum and equal to 2 on the scale adopted. The first
curve proceeding outwards is the locus of points at which w=1. The next
is the nodal line, separating the regions of opposite displacement. The
remaining curves taken in order give the displacements - 1,- 2, -3. The
numerically greatest negative displacement occurs at the corners of the
square, where it amounts to 2 x 1645 = 3'290."

Dangling\footnote{ On the nodal lines of a square plate. Phil. Mag. August, 1873.}

The nodal curve thus constructed agrees pretty closely with the
observations of Strehlke
\footnote{Pogg. Ann. Vol. CXLVI p. 319}. His
results, which refer to three carefully worked plates of glass, are
embodied in the following polar equations:

40143 0171) 00127 r= 40143 + 0172) (4t) + 00127 (8t) 4019 0168) 0013

Fig. 42.

the centre of the square being pole. From these we obtain for the radius
vector parallel to the sides of the square (t=0) 0.41980, 0.41981,
0.4200, while the calculated result is 0.4154. The radius vector
measured along a diagonal is 3856, 3855, 3864, and by calculation 3900.

By crossing the network in the other direction we obtain the locus of
points for which :

u1(x/a)-u1(y/a)

is constant, which are the curves of constant displacement for that mode
in which the diagonals are nodal. The pitch of the vibration is
(according to theory) the same in both cases.

Fig. 43.

The primitive modes represented by w=, or w=4, (

may be combined in like manner, Fig. 43 shews the nodal curve

for the vibration w = u2 (x/a) pm + u2(y/a) (5).

The form of the curve is the same relatively to the other diagonal, if
the sign of the ambiguity be altered.

227. The method of superposition does not depend for its application on
any particular form of normal function. Whatever the form may be, the
mode of vibration, which when p= 0 passes into that just discussed, must
have the same period, whether the approximately straight nodal lines are
parallel to az or to y. If the two synchronous vibrations be superposed,
the resultant has still the same period, and the general course of its
nodal system may be traced by means of the consideration that no point
of the plate can be nodal at which the primitive vibrations have the
same sign. To determine exactly the line of compensation, a complete
knowledge of the primitive normal functions, and not merely of the
points at which they vanish, would in general be necessary. Doctor Young
and the brothers Weber appear to have had the idea of superposition as
capable of giving rise to new varieties of vibration, but it is to Sir
Charles Wheatstone' that we owe the first systematic application of it
to the explanation of Chladni's figures. The results actually obtained
by Wheatstone are however only very roughly applicable to a plate, in
consequence of the form of normal function implicitly assumed. In place
of Fig. 42 (itself, be it remembered, only an approximation) Wheatstone
finds for the node of the compound vibration the inscribed square shewn
in Fig. 44,

O

This form is really applicable, not to a plate vibrating in virtue of
rigidity, but to a stretched membrane, so supported that every point of
the circumference is free to move along lines perpendicular to the plane
of the membrane. The boundary condition applicable under these
circumstances is m0; and it is easy to shew that the normal functions
which involve only one coordinate are w =cos (m ), or w = (m w), the
origin being

at a corner of the square. Thus the vibration

wy = 008 + (27Y st (1) a a

Dangling\footnote{ Phil. Trans. 18338,}

has its nodes determined by

whence +y=4a or ga, or e--y=+ 42, equations which represent the
inscribed square. :

If w = 008 Te - cog 27Y (2),

the nodal system is composed of the two diagonals. This result, which
depends only on the symmetry of the normal functions, is strictly
applicable to a square plate. When m= 3, 8ara 3iry

W = () + () (3),

and the equations of the nodal lines are

5 aty=5, a, 3 -y=tR

shewn in Fig. 45, If the other sign be taken, we obtain a similar figure
with reference to the other diagonal.

When m = 4,

With the other sign

we obtain aty=5, a, 2, x-y=0, 5 Fig. 47),

representing a system composed of the diagonals, together with the
inscribed square. .

These forms, which are strictly applicable to the membrane, resemble the
figures obtained by means of sand on a square plate more closely than
might have been expected. The sequence of tones is however quite
different. From Article 176 we see that, if were zero, the interval
between the form (43) derived from three primitive nodes, and (41) or
(42) derived from two, would be 1:4629 octave; and the interval between
(41) or (42) and (46) or (47) would be 24358 octaves. Whatever may be
the value of w the forms (41) and (42) should have exactly the same
pitch, and the same should be true of (46) and (47). With respect to the
first mentioned pair this result is not in agreement with Chladni's
observations, who found a difference of more than a whole tone, (42)
giving the higher pitch. If however (42) be left out of account, the
comparison is more satisfactory. According to theory (w=0), if (41) gave
d, (43) should give g' -, and (46), (47) should give g''+. Chladni found
for (48) g'+, and for (46), (47) y and Vk + respectively,

228. The gravest mode of a square plate has yet to be considered. The
nodes in this case are the two lines drawn through the middle points of
opposite sides. That there must be such a mode will be shewn presently
from considerations of symmetry, but neither the form of the normal
function, nor the pitch, has yet been determined, even for the
particular case of =0. A rough calculation however may be founded on an
assumed type of vibration.

If we take the nodal lines for axes, the form w= ay satisfies v w =0, as
well as the boundary conditions proper for a free edge at all points of
the perimeter except the actual corners. This is in fact the form which
the plate would assume if held at rest by four forces numerically equal,
acting at the corners perpendicularly to the plane of the plate, those
at the ends of one diagonal being in one direction, and those at the
ends of the other diagonal in the opposite direction. From this it
follows that w= xy would be a possible mode of vibration, if the
mass of the plate were concentrated equally in the four corners. By (3)
Article 214, we see that

inasmuch as dat = dy s = da dy Pt

For the kinetic energy, if p be the volume density, and M the additional
mass at each corner,

T= 4p" (pt) f- : i dx dy + Ma

= ip (pt) 9+ ul (2).

Hence (3),

where M' denotes the mass of the plate without the loads. This result
tends to become accurate when WM is relatively great; otherwise by
Article 89 it is sensibly less than the truth. But even when M=0, the
error is probably not very great. In this case we should have

giving a pitch which is somewhat too high. The gravest mode next after
this is when the diagonals are nodes, of which the pitch,

if = 0, would be given by (5), (see Article 174).

We may conclude that if the material of the plate were such that w= 0,
the interval between the two gravest tones would be somewhat greater
than that expressed by the ratio 1318. Chladni makes the interval a
fifth.

229. That there must exist modes of vibration in which the two shortest
diameters are nodes may be inferred from such considerations as the
following. In Fig. (48) suppose that GH
is a plate of which the edges HO; GO are supported, and the edges GO, CH
free. This plate, since it tends to a definite position of equilibrium,
must be capable of vibrating in certain fundamental modes. Fixing our
attention on one of these, let us conceive a distribution of w over the
three remaining quadrants, such that in any two that adjoin, the values
of w are equal and opposite at points which are the images of each other
in the line of separation. If the whole plate vibrate according to the
law thus determined, no constraint will be required in order to keep the
lines GE, FH fixed, and therefore the whole plate may be regarded as
free. The same argument may be used to prove that modes exist in which
the diagonals are nodes, or in which both the diagonals and the
diameters just considered are together nodal. The principle of symmetry
may also be applied to other forms of plate. We might thus infer the
possibility of nodal diameters in circle, or of nodal principal axes in
the ellipse: When the
boundary is a regular hexagon, it is easy to see that Figs. (49), (50),
(51) represent possible forms.

Fig. 49.

Fig. 51.

It is interesting to trace the continuity of Chladni's figures, as the
form of the plate is gradually altered. In the circle, for example, when
there are two perpendicular nodal diameters, it is a matter of
indifference as respects the pitch and the type of vibration, in what
position they be taken. As the circle develops into a square by throwing
out corners, the position of these diameters becomes definite. In the
two alternatives the pitch of the vibration is different, for the
projecting corners have not the same efficiency in the two cases. The
vibration of a square plate shewn in Fig. (42) corresponds to that of a
circle when there is one circular node. The correspondence of the graver
modes of a hexagon or an ellipse with those of a circle may be traced in
like manner.

230. For plates of uniform material and thickness and of invariable
shape, the period of the vibration in any fundamental mode varies as the
square of the linear dimension, provided of course that the boundary
conditions are the same in all the cases compared. When the edges are
clamped, we may go further and assert that the removal of any external
portion is attended by a rise of pitch, whether the material and the
thickness be uniform, or not.

Fig. 62.

Let AB be a part of a clamped edge
(it is of no consequence whether the remainder of the boundary be clamped, or net), and
let the piece ACBD be removed, the new edge ADB being also clamped. The
pitch of any fundamental vibration is sharper than before the change.
This is evident, since the altered vibrations might be obtained from the
original system by the introduction of a constraint clamping the edge
ADB. The effect of the constraint is to raise the pitch of every
component, and the portion ACBD being plane and at rest throughout the
motion, may be removed. In order to follow the sequence of changes with
greater security from error, it is best to suppose the line of clamping
to advance by stages between the two positions ACB, ADB. For example,
the pitch of a uniform clamped plate
in the form of a regular hexagon is lower than for the inscribed circle
and higher than for the circumscribed circle.

When a plate is free, it is not true that an addition to the edge always
increases the period. In proof of this it may be sufficient to notice
a particular case.

AB is a narrow thin plate, itself without inertia but carrying loads at
A, B,C. It is clear that the addition to the breadth
indicated by the dotted line would augment the stiffness of the bar, and
therefore lessen the period of vibration. The same consideration shews
that for a uniform free plate of given area there is no lower limit of
pitch; for by a sufficient elongation the period of the gravest
component may be made to exceed any assignable quantity. When the edges
are clamped, the form of gravest pitch is doubtless the circle.

If all the dimensions of a plate, including the thickness, be altered in
the same proportion, the period is proportional to the linear dimension,
as in every case of a solid body vibrating in virtue of its own
elasticity.

The period also varies inversely as the square root of Young's modulus,
if be constant, and directly as the square root of the mass of unit of
volume of the substance.

231. Experimenting with square plates of thin wood whose - grain ran
parallel to one pair of sides, Wheatstone' found that the pitch of the
vibrations was different according as the approximately straight nodal
lines were parallel or perpendicular to the fibre of the wood. This
effect depends on a variation in the flexural rigidity in the two
directions. The two sets of vibrations having different periods cannot
be combined in the usual manner, and consequently it is not possible to
make such a plate of wood vibrate with nodal diagonals. The inequality
of periods may however be obviated by altering the ratio of the sides,
and then the ordinary mode of superposition giving nodal diagonals is
again possible. This was verified by Wheatstone.

Dangling\footnote{ Phil. Trans. 1833}

A further application of the principle of superposition is due to Koenig
\footnote{Pogg Ann. 1864, CXXII. p. 238}. In order that two modes of
vibration may combine, it is only necessary that the periods agree. Now
it is evident that the sides of a rectangular plate may be taken in such
a ratio, that (for instance) the vibration with two nodes parallel to
one pair of sides may agree in pitch with the vibration having three
nodes parallel to the other pair of sides. In such a case new nodal
figures arise by composition of the two primary modes - of vibration, .

232. When the plate whose vibrations 2 are to be considered is naturally
curved, the difficulties of the question are generally much increased.
But there is one case in which the complication due to curvature is more
than compensated by the absence of a free edge; and this case happens to
be of considerable interest, as being the best representative of a bell
which at present admits of analytical treatment.

A long cylindrical shell of circular section and uniform thickness is
evidently capable of vibrations of a flexural character in which the
axis remains at rest and the surface cylindrical, while the motion of
every part is perpendicular to the generating lines. The problem may
thus be treated as one of two dimensions only, and depends upon the
consideration of the potential and kinetic energies of the various
deformations of which the section is capable. The same analysis also
applies to the corresponding vibrations of a ring, formed by the
revolution of a small closed area about an external axis.

The cylinder, or ring, is susceptible of two classes of vibrations
depending respectively on extensibility and flexural rigidity, and
analogous to the longitudinal and lateral vibrations of straight bars.
When, however, the cylinder is thin, the forces resisting bending become
small in comparison with those by which extension is opposed; and, as in
the case of straight bars, the vibrations depending on bending are
graver and more important than those which have their origin in
longitudinal rigidity. In the limiting case of an infinitely thin shell
(or ring), the flexural vibrations become independent of any extension
of the circumference as a whole, and may be calculated on the
supposition that each part of the circumference retains its natural
length throughout the motion.

But although the vibrations about to be considered are analogous to the
transverse vibrations of straight bars in respect of depending on the
resistance to flexure, we must not fall into the common mistake of
supposing that they are exclusively normal. It is indeed easy to see
that a motion of a cylinder or ring in which each particle is displaced
in the direction of the radius would be incompatible with the condition
of no extension. In order to satisfy this condition it is necessary to
ascribe to each part of the circumference a tangential as well as a
normal motion, whose relative magnitudes must satisfy a certain
differential equation. Our first step will be the investigation of this
equation.

233. The original radius of the circle being a, let the equilibrium
position of any element of the circumference be defined by the vectorial
angle 6. During the motion let the polar coordinates of the element
become

r=atr, 6=0+400.

If ds represent the arc of the deformed curve corresponding to ad , we
have

(ds) = (ad6) = (ddr) + 7 (d6 + d80);

whence we find, by neglecting the squares of the small quantities br,
50,

as the required relation.

In whatever manner the original circle may be deformed at time c, 5r may
be expanded by Fourier's theorem in the series

r=a A (0)+ B (0) + A (20) + B (26)+... +A (n) + B (n) + ...(2),

and the corresponding tangential displacement required by the condition
of no extension will be

= - A, (0)+ B (0+... - 4s (nd) +22 (28 = (3);

the constant that might be added to 56 being omitted.

If cad denote the mass of the element ad , the kinetic energy T' of the
whole motion will be

= +3 (dp + B+...

+ (1 + =) (42453) + os

the products of the coordinates A,, 8, disappearing in the integration.

We have now to calculate the form of the potential energy V. Let p be
the radius of curvature of any element ds, then for the

corresponding element of V we may take 4Bds (3), where B is a
constant depending on the material and on the thickness. Thus

V=4Ba i (5) dO (5). Now Ii 4, Pp dg" and

11 u=r=cil- A()- B (g) -..., for in the small terms the distinction
between c and may be neglected. Hence

5 = 3 S (nt 1)(4, (nd), 1 p and

v-2 I = (n - 1) (A (nd + B, (n6) dd

7 By (n = 1)"(A2 + BY) , (6)

in which the summation extends to all positive integral values of n.

The term for which n= 1 contributes nothing to the potential energy, as
it corresponds to a displacement of the circle as a whole, without
deformation.

We see that when the configuration of the system is defined as above by
the coordinates A,, B,, etc., the expressions for Z and V involve only
squares; in other words, these are the normal coordinates, whose
independent harmonic variation expresses the vibration of the system.

If we consider only the terms involving (0), (n) , we have by taking the
origin of suitably,

Br =A (n)8, 80=- Ae (RB (1),

while the equation defining the dependence of A, upon the time is .

1 y B oa (1 +=) Alt S (n= 1A, =0 se (8), from which we conclude that, if
A, varies as (pt-e), B n'(n -1) p= cat (9).

This result was given by Hoppe for a ring in a memoir published in
Crelle, Bd. 63,1871. His method, though more complete than the
preceding, is less simple, in consequence of his not recognising
explicitly that the motion contemplated corresponds to complete
inextensibility of the circumference.

second edition, page 399, pdfpage 385

[In the application of (9) to a ring we have, § 192 a,
\[ \frac{B}{\sigma}= \frac{c^2}{4} \frac{q}{\rho}\]
(10).  where q is Youngs modulus, \(\rho\) the volume density, and c the
radius of the circular section. For the cylindrical shell.
18, § 235 g,
\[ \frac{B}{\sigma}= \frac{4 m h^2}{3(m+n)\rho}\]
(11)
\(2h\) denoting the thickness, and m, n the elastic constants in
Thomson and Tail’s notation.]

According to Chladni the frequencies of the tones of a ring are as

If we refer each tone to the gravest of the series, we find for the
ratios characteristic of the intervals

2.778, 5.445, 9, 13.44, etc,

The corresponding numbers obtained from the above theoretical formulae,
by making n successively equal to 2, 3, 4, etc., are

2.828, 5.423, 8.771, 12.87, etc.,
agreeing pretty nearly with those found experimentally.

[Observations upon the tones of thin metallic cylinders, open
at one end, have been made by Fenknor\footnote{Wied. Ann. vol. 8, p 185, 1879.}
Since the pitch proves to be very nearly independent of the height of the
cylinders, the vibrations may be regarded as approximately two-dimensional.  In
accordance with (9), (11), Fenknor found the frequency proportional to the
thickness directly, and the square of the radius inversely. As regards the
sequence of tones from a given cylinder\footnote{Melde, Akustik, Leipzig, 1883.
p.223 }, the numbers, Deferred to the gravest (s=2) as unity,
were 2.67, 5.00, 8.00, 12.00, etc. The agreement with (9) would
be improved if these numbers were raised by about \(1/12\) part
equivalent to an alteration in the pitch of the gravest tone.

The influence of rotation of the shell about its axis hii' boon
examined by Bryan\footnote{Proc. Climb. Phil. Soc, vol. VII. p. 101, 1890}
It appears that the nodes are carried round,
but with an angular velocity less than that of the rotation
If the latter be denoted by \(\omega\), the nodal angular velocity is
\[ \frac{ s^2-1 }{ s^2+1 } \omega \]
]


234. When n=1, the frequency is zero, as might have been anticipated.
The principal mode of vibration corresponds to n = 2, and has four
nodes, distant from each other by 90. These so called nodes are not,
however, places of absolute rest, for the tangential motion is there a maximum.
In fact the tangential vibration at these points is half the maximum normal
motion, In general for the n term the maximum tangential motion is of the
maximum normal motion, and occurs at the nodes of the latter.

When a bell-shaped body is sounded by a blow, the point of application
of the blow is a place of maximum normal motion of the resulting
vibrations, and the same is true when the vibrations are excited by a
violin-bow, as generally in lecture room experiments. Bells of glass,
such as finger-glasses, are however more easily thrown into regular
vibration by friction with the wetted finger carried round the
circumference. The pitch of the resulting sound is the same as of that
elicited by a tap with the soft part of the finger; but inasmuch as the
tangential motion of a vibrating bell has been very generally ignored,
the production of sound in this manner has been felt as a difficulty. It
is now scarcely necessary to point out that the effect of the friction
is in the first instance to excite tangential motion, and that the point
of application of the friction is the place where the tangential motion
is greatest, and therefore where the normal motion vanishes.

235. The existence of tangential vibration in the case of a bell was
verified in the following manner. A so-called air-pump receiver was
securely fastened to a table, open end uppermost, and set into vibration
with the moistened finger. A small chip in the rim, reflecting the light
of a candle, gave a bright spot whose motion could be observed with a
Coddington lens suitably fixed. As the finger was carried round, the
line of vibration was seen to revolve with an angular velocity double
that of the finger; and the amount of excursion (indicated by the length
of the line of light), though variable, was finite in every position.
There was, however, some difficulty in observing the correspondence
between the momentary direction of vibration and the situation of the
point of excitement. To effect this satisfactorily it was found
necessary to apply the friction in the neighbourhood of one point. It
then became evident that the spot moved tangentially when the bell was
excited at points distant therefrom 0, 90, 180, or 270 degrees; and
normally when the friction was applied at the intermediate points
corresponding to 45, 135, 225 and 315 degrees. Care is sometimes
required in order to make the bell vibrate in its gravest mode without
sensible admixture of overtones.

If there be a small load at any point of the circumference, a slight
augmentation of period ensues, which is different according as the
loaded point coincides with a node of the normal or of the tangential
motion, being greater in the latter case than in the former. The sound
produced depends therefore on the place of excitation; in general both
tones are heard, and by interference give rise to beats, whose frequency
is equal to the difference between the frequencies of the two tones.
This phenomenon may often be observed in the case of large bells.


\section{235 a Church Bells}
In determining the number of nodal meridians (2s) corresponding to any
particular tone of a bell, advantage may be taken of beats, whether due to
accidental irregularities or introduced for the purpose by special loading
(compare §§ 208, 209). By tapping cautiously round a circle of latitude the
places may be investigated where the beats disappear, owing to the absence of
one or other of the component tones. But here a decision must not be made too
hastily. The inaudibility of the beats may be favoured by an unsuitable
position of the ear or of the mouth of the resonator used in connection with
the ear. By travelling round, a situation is soon found where the observation
can be made to the best advantage. In the neighbourhood of the place where the
blow is being tried there is a loop of the vibration which is most excited and
a (coincident) node of the vibration which is least excited. When the ear is
opposite to a node of the first vibration, and therefore to a loop of the
second, the original inequality is redressed, and distinct beats may be heard
even though the deviation of the blow from a nodal point may be very small. Th
e accurate determination in this way of two consecutive places where no beats
are generated is all that is absolutely necessary for the purpose in view. The
ratio of the entire circumference of the circle of latitude to the arc between
the pints in question is in fact 4s.  Thus, if the arc between consecutive points
proved to be \(45^{\circ}\), we should infer that we were dealing with the case
of \(s=2\), in which the deformation is elliptical. As a greater security
against error, it is advisable in practice to determine a larger number of points
where no beats occur. Unless the deviation from symmetry be considerable, these
points should be uniformly distributed along the circle of latitude\footnote{The
bells, or gongs, as they are sometimes called, of striking clocks often give
disagreeable beats. A remedy may be found in a suitable rotation of the bell
round its axis.}

In the above process for determining nodes we are supposed to hear distinctly
the tone corresponding to the vibration under investigation. For this purpose
the beats are of assistance in directing the attention', but in dealing with
the more difficult subjects, such as church bells, it is advisable to have
recourse to resonators. A set of v, Helmholtz s pattern, as manufactured by
Koenig, are very convenient. The one next higher in pitch to the tone under
examination is chosen and tuned by advancing the finger across the aperture.
Without the security afforded by resonators, the determination of the octave is
very uncertain.

The only class of bells, for which an approximate theory can be given, are
those with thin walls, 233, 235 c. Of such the following glass balls may be
regarded as examples:
\todo[inline]{smudged table, sharps or flats}
I. c',e'',c''', II. a, c'',b''  III. ?, b''
The value of s for the gravest tone was 2, for the second 3, and for the third
tone 4.

Similar observations have been made upon a so-called hemispherical bell, of
nearly uniform thickness, and weighing about 3 cwt.
\todo[inline]{smudged table, ewt} Four tones could be plainly
heard,
 e, f', e'', b''
the pitch being taken from a harmonium. The gravest tone has a long duration.
When the bell is struck by a hard body, the higher tones are at first
predominant, but after a time they die away, and leave in possession of the
field. If the striking body be soft, the original preponderance of the higher
elements is less marked.

By the method described there was no difficulty in shewing that the four tones
correspond respectively to = 2, 3, 4, 5. Thus for the gravest tone the
vibration is elliptical with 4 nodal meridians, for the next tone there are
nodal meridians, and so on.

Tapping along a meridian shewed that the sounds became less clear as the edge
was departed from, and this in a continuous manner with no suggestion of a
nodal circle of latitude. A quest ion to which we shall recur in connection
with church bells here suggests itself. Which of the various coexisting tones
characterizes the pitch of the bell as a whole  It would appear to be the third
in order, for the founders gave the pitch as E natural.

In church bells there is great concentration of metal at the ‘sound-bow” where
the clapper strikes, indeed to such an extent that we can hardly expect much
correspondence with what occurs in the case of thin uniform bells. But the
method already described suffices to determine the number of nodal meridians
for all the more important tones. From a bell of cwt. bv Mears and Stainbank
(j tones could be obtained, viz. :
c, c, +. , (r (4) 4) n) (6) (S)

The pitch of this bell as given by the makers is d' so that it is the fifth in
the above series of tonUsS which characterizes the bell. The number of nodal
meridians in the various components is indicated within the parentheses. Thus
in the case af the lone c' there are 4 nodal meridians. A similar method of
examination along a meridian shewed that there was no nodal circle of latitude.
At the same time differences of intensity were observed. This tone is most
fully developed when the blow is delivered about midway between the crown and
the rim of the bell.

The next tone is c”. Observation shewed that for this vibration also there are
four, and but four, nodal meridians. But now there is a well-defined nodal
circle of latitude, situated about a quarter of the way up from the rim towards
the crown. As heard with a resonator, this tone disappears when the blow is
accurately delivered at some point of this circle, but revives with a very
small displacement on either side. The nodal circle and the four meridians
divide the surface into segments, over each of which the normal motion is of
one sign.

To the tone \(f''\) corresponds 6 nodal meridians. There is no well-defined nodal
circle. The sound Ls indeed very faint when the tap is much displaced from the
sound-bow; it was thought to fall to a minimum when a position about half-way
up was reached.

The three graver tones are heard loudly from the sound-bow.  But the next in
order, 6l7, is there scarcely audible, unless the blow is delivered to the rim
itself in a tangential direction. The maximum effect occurs about half-way up.
Tapping round the circle revealed G nodal meridians.

The fifth tone, d', is heard loudly from the sound-bow, but soon falls off when
the locality of the blow is varied, and in the upper three-fourths of the bell
it is very faint. No distinct circular node could be detected. Tapping round
the circumference shewed that there were 8 nodal meridians.

The highest tone recorded, was not easy of observation, and the mode of
vibration could not be fixed satisfactorily.

Similar results have been obtained from a bell of 4 cwt., cast
by Taylor of Loughborough for Ampton church. The nominal
pitch (without regard to octave) was d, and the following were the
tones observed:
\begin{tabular}{ccccc}
e'-2& d''-6& f''+4& b''-b& d''', g''' \\
(4) & (4)  &(6)   &(6)   &(8) \\
\end{tabular}
In the specification of pitch the numerals following the note indicate by how
much the frequency for the bell differed from that of the harmonium employed as
a standard. Thus the gravest tone e'b gave 2 beats per second, and was fiat.
When the number exceeds 3, it is the result of somewhat rough estimation, and
cannot be trusted to be quite accurate. Moreover, as has been explained, there
care in strictness two frequencies under each head, and these often differ
sensibly. In the case of the 4th tone, 6 1 - h means that, as nearly as could
be judged, the pitch of the bell was midway between the two specified notes of
the harmonium.

Observations in the laboratory upon the above-mentioned bells having settled
the modes of vibration corresponding to the five gravest tones, other bells of
the church pattern could be sufficiently investigated by simple determinations
of pitch. The results are collected in the following table
\footnote{On Bells, Phil. Mag., vol. 29. p. 1. 1890} and include,
besides those already given, observations upon a Belgian bell, the property of
Mr Haweis, and upon the five bells of the Terling peal. As regards
the nominal pitch of the latter bells,
several observers concurred in fixing the notes of the peal as
\[ f_{\sharp}, g_{\sharp}, a_{\sharp}, b, c_{\sharp} \]
no attention being paid to the question of the octave.

page 393,  pdf page 407, table, Actual Pitch by Harmonium, Pitch referred to fifth

Examination of the table reveals the remarkable fact that in every case of
the English bells it is the 5th tone in order which agrees with the nominal
pitch, and that, with the exception of Terling (4), no other tone shews such
agreement \footnote{In this comparison the gravest tone is disregarded.}.
Moreover, as appeared most clearly in the case of the bell cast by
Mears and Stainbank, the nominal pitch, as given by the makers, is an octave
below the only corresponding tone.

The highly composite, and often discordant, character of the sounds of bells
tends to explain the discrepancies sometimes, manifested in estimations of
pitch. Mr Simpson, who has devoted much attention to the subject, has put
forward strong arguments for the opinion that the Belgian makers determine the
pitch of their bells by the tone 2nd in order in the above series, so that for
instance the pitch of Terling (3) would be a and not at. In subordination to
this tone they pay attention also to the next (the 3rd in order), classifying
their bells according to the character
of the third, whether
major or minor, so compounded. Thus in Terling (3) the interval, a to c , is a
major third. The comparative neglect with which the Belgians treat the 5th
tone, regarded almost exclusively by English makers, may perhaps be explained
by a less prominent development of this tone in Belgian bells, and by a
difference in treatment. When a bell is sounded alone, or with other bells in
a comparatively slow succession, attention is likely to concentrate itself upon
the graver and more persistent elements of the sound rather than upon the
acuter and more evanescent elements, while the contrary may be expected to
occur when bells follow one another rapidly ii a peal.

In any case the false octaves with which the Table abounds are simple facts of
observation, and wv may well believe that their correction would improve the
general effect. Especially should the octave between the 2nd tone and the 5th
tone be made true.  Probably the lower octave of the gravest, or hum-note, as
it is called by English founders, is of less importance. The same may be said
of the fifth, given by the 4th tone of the series, which is much less
prominent. The variations recorded in the Table would seem to shew that no
insuperable obstacle stands in the way of obtaining accurate harmonic relations
among the various tones.

No immediate explanation has been given of the form adopted for church bells.
It appears both from experiment and from the theory of thin shells that this
form is especially stiff, as regards the principal mode of deformation (s = 2),
to forces applied normally and near the rim. Possibly the advantage of this form
lies in its rendering less prominent the gravest component of the sound, or the
hum-note.



\section{My notes on Rayleigh's description of plates}
What is the strain energy of bending in the theory of elastic plates?

Theorem.

isotropic elastic material, E, nu, L, G,...
infinitesimal strain, e, plane stress,
eo := e(1:2,1:2),
\[
to := \trace(eo),
\]
  do := |eo|,.
  U = (1/2)  E /(1-nu nu)  int  ( to to - 2 (1-nu) do ) dA

Lemmas.
A.
\[
{\tr(H)}^2 = H:H + 2 |H|  \rightarrow  to to = eo : eo + 2 do
\]

B.  (L + 2G)/(L+G) =  2 (1-nu)
C.  2G = E /(1+nu)

Proof of Theorem.  U = 1/2 int s:e dA.
so = s(1:2,1:2),  s = diag(so,0), e = diag(eo,e330),

\[
   \sigma = \lambda \tr (e) I + 2 G  e
\]
\begin{verbatim}
   so = L trace(e) I + 2 G eo
to = I : eo,  L trace(e) + 2 G e33 = 0,  e33 =  - to  L/(L+2 G)
trace(e) = to + e33 = to ( 1 - L/(L+2G)) = to  2G/(L+2G)
so =  to (2 LG)/(L+2G)  I + 2 G eo
\end{verbatim}

\begin{verbatim}
s:e = so : eo = to to  2 LG/(L+2G) +  2G  (eo:eo) +
     2 G ( to to  L/(L+2G) +    eo:eo  )
     2 G ( to to  L/(L+2G) +    to to - 2 do   )
     4 G ( to to  (L+G)/(L+2G)  -  do   )
    = 4 G (L+G)/(L+2G)  ( to to - (L+2G)/(L+G) do )
  = 2 G /(1-nu)   ( to to - 2 (1-nu)(do )
  = E /(1-nu nu)   ( to to - 2 (1-nu)(do )
end of proof.
\end{verbatim}


Kirchhoff - Love,  eo(x,y,z) = - z Hess(w),
\[
\int_{-h/2}{h/2} z^2 dz = (2/3) (h/2)^3 = h^3/12
\]
Something like that.  Sometimes  [-h,h],

J. Strutt explains in terms of
Hess(w),  eigenvalues  1/p1 and 1/p2,
\[
U = c ( 1/p1^2 + 1/p2^2 + f /p1 p2 ),  0 < f < 1.
\]
That is consistent.
\printbibliography[heading=bibintoc,title={References}]
\end{document}
